% mage-auto-heading: the-two-problems-in-software
% mage-auto-heading: the-methodology-extends-beyond-software
% Authoritative LaTeX manuscript; Markdown is a generated compatibility projection.

% <!-- part-title: The New Engineering Problem -->

% <!-- chapter-title: Two Problems the Method Must Solve -->

\chapter{Two Problems the Method Must Solve}

\section{The Two Problems in Software}\label{heading:the-two-problems-in-software}

% <!-- point: substrate-leaves-two-problems | The engineered environment can improve the agent’s implementation success and support engineers in choosing requirements and establishing whether the result meets them; these aims lead to the Representation Problem and the Enforcement Problem. | terms: engineering-substrate, foundation-model, agentic-harness -->

The probability model in \magevalue{sec:reasoning-horizon} gives us one aim: increase the agent’s chance of implementing all the required features correctly.\footnote{Or at least, correctly enough that we can finish the job economically. Human code review and cleanup, or a final pass with a more capable but more expensive model, may supply what the initial agent leaves undone. The relevant cost is that of the whole process, including finding and fixing mistakes; the result must still meet its requirements.} The engineered environment—the structures and practices we arrange around the agent’s work—can also support engineers in deciding what those requirements should be and establishing whether the result meets them.

Better context can help an agent make a consequential judgment correctly. A useful model can sometimes change the task itself: it makes the relevant property explicit enough to analyze or check without reconstructing the implementation. The first improves the agent’s chances of getting the answer right; the second can reduce how often the question must be answered through inference.

These two aims—the agent’s implementation success, and support for engineers choosing requirements and establishing whether the result meets them—lead to two related but distinct engineering problems:

% <!-- point: the-two-problems-representation-and-enforcement | The two problems, related but distinct: the Representation Problem (how to represent a large system so relevant properties are reasoned about without reconstruction) answered by Modeling in Chapter 2, and the Enforcement Problem (how intent comes to bind autonomous work so satisfying obligations does not depend on the producer self-certifying) answered by Alignment in Chapter 3. | terms: engineering-substrate, modeling-principle, alignment-principle, reasoning-horizon, injection-point -->

% <!-- index-def: representation-problem -->

% <!-- index-def: enforcement-problem -->

\begin{enumerate}
\def\labelenumi{\arabic{enumi}.}
\tightlist
\item
  \textbf{The Representation Problem.} Large systems contain far more detail than any one engineering question needs, while the active state available to a reasoner is finite. How should a large, evolving system be represented so that the properties relevant to an engineering question can be reasoned about without continually reconstructing them from implementation? \magevalue{chapter:modeling} develops the corresponding part of the MAGE methodology: Modeling.
\item
  \textbf{The Enforcement Problem.} Autonomous work can be constrained and checked at the interfaces through which it acts, so that the agent is not its own arbiter of validity. How should engineering intent be enforced over autonomous work, so that satisfying important obligations does not depend on the agent remembering, agreeing, or self-certifying each time? \magevalue{chapter:alignment} develops the corresponding part of the MAGE methodology: Alignment.
\end{enumerate}

% <!-- point: modeling-and-alignment-complement-one-another | Modeling represents selected system properties; Alignment constrains work and evaluates results. In DocAble, the content representation supports comparison and the validator fails a job that loses original extracted content. | terms: modeling-principle, alignment-principle -->

The two activities complement one another. Modeling makes selected system properties explicit enough to reason about; Alignment uses represented obligations to constrain work and evaluate its results. In the DocAble example, the content representation lets a validator compare the original and revised documents. The validator fails the job if the revised document is missing any of the original extracted content. Together, these structures carry an engineering decision forward into later work. \autoref{substrate-to-method} connects the two problems to the method.

\begin{figure}[htbp]
\centering
\includegraphics[width=\linewidth]{book/latex/assets/substrate-to-method-7fa75819e3.pdf}
\caption{\emph{From substrate to method.} Finite reasoning state creates the Representation Problem; probabilistic autonomous work creates the Enforcement Problem. Chapters 2 and 3 develop MAGE\textquotesingle s corresponding answers: Modeling and Alignment.}
\label{substrate-to-method}
\end{figure}

\section{The Methodology Extends Beyond Software}\label{heading:the-methodology-extends-beyond-software}

% <!-- point: the-problem-is-larger-than-software | Delegating work to agents raises the same engineering questions across domains, but the useful representations and credible evidence differ. -->

The same questions will arise wherever humans delegate work to agents. However, the useful representations and credible evidence will differ by domain. Software engineering has architectures, dependency graphs, and tests. Mechanical engineering has geometry, tolerances, simulations, and physical tests. Accounting has ledgers, reconciliations, and audit evidence.

The balance between delegation, explicit representation, and enforcement will change as agents improve. MAGE does not prescribe that balance in advance. It asks engineers to decide which requirements to represent and enforce, and which choices to leave open as degrees of freedom for the agent.

This book develops MAGE through software engineering, drawing on the DocAble case study. \magevalue{chapter:modeling} asks how to represent the knowledge needed for engineering decisions. \magevalue{chapter:alignment} asks how selected obligations can govern the work. \magevalue{chapter:theory} develops the theory and examines its scope and limits. \magevalue{chapter:profession} considers its implications beyond software.
